\subsubsection{Transversal Gates}

In fault-tolerant quantum computing, one of the \hl{main objectives is to apply logical quantum gates without allowing physical errors to spread uncontrollably through the encoded qubits.} Let's suppose one logical qubit\footnote{A logical qubit is an encoded qubit constructed from many physical qubits in order to protect quantum information against errors. Instead of storing the quantum state in a single fragile qubit, the information is distributed redundantly across multiple physical qubits.} is encoded into many physical qubits:
\begin{equation*}
    \text{Logical Qubit} \rightarrow q_1, q_2, \ldots, q_7
\end{equation*}
It means that a single logical qubit is represented by a collection of physical qubits. Now imagine, for example, that $q_3$ contains an error, and all the other qubits are fine. This situation is still manageable because we can detect and correct the error in $q_3$ without affecting the other qubits.

\highspace
\textcolor{Red2}{\faIcon{exclamation-triangle} \textbf{The Dangerous Situation.}} Now suppose we apply a complicated gate that strongly couples all qubits together. Then, if $q_3$ has an error, it can spread to all the other qubits, making the entire logical qubit corrupted and uncorrectable.

\highspace
\textcolor{Green3}{\faIcon{check-circle} \textbf{The core goal of Fault Tolerance.}} So the main objective becomes: if \textbf{one physical qubit fails, the failure should remain localized and not spread to other qubits}. \hl{This is where transversal gates come into play.}

\highspace
\begin{definitionbox}[: Transversal Gates]
    A \definition{Transversal Gate} is a logical quantum \textbf{gate} implemented by \textbf{applying independent operations} to corresponding \textbf{physical qubits of an encoded logical qubit}.

    \highspace
    Suppose a logical qubit is encoded into $n$ physical qubits. A transversal implementation applies the same unitary operation independently to each physical qubit:
    \begin{equation*}
        U_L = U^{\otimes n}
    \end{equation*}
    Where:
    \begin{itemize}
        \item $U$ is the physical gate;
        \item $U_L$ is the resulting logical gate;
        \item $n$ is the number of physical qubits used in the encoding.
    \end{itemize}
\end{definitionbox}

\highspace
In simple terms, a transversal gate applies operations independently to each physical qubit. So no qubit interacts with another qubit inside the same block. Therefore if $q_3$ was wrong before, only $q_3$ will be affected by the gate, and \textbf{the error will \emph{not} spread to the other qubits}.

\newpage

\begin{flushleft}
    \textcolor{Green3}{\faIcon{book} \textbf{Error Containment and Fault Tolerance}}
\end{flushleft}
The \textbf{most important property} of transversal gates is that \hl{they prevent error propagation inside the encoded block.}

\highspace
Suppose one physical qubit contains an error before the gate is applied. Since each qubit is manipulated independently, the error remains localized and does not spread to many other qubits.

\highspace
Furthermore, transversal gates are \textbf{inherently fault-tolerant}. In general, a fault-tolerant operation must ensure that:
\begin{itemize}
    \item A small physical error remains small;
    \item A single faulty qubit does not corrupt the entire logical qubit.
\end{itemize}
Transversal gates naturally satisfy this requirement \hl{because there are no interactions between qubits within the same code block during the gate operation.} Therefore, even if one qubit experiences an error, it does not affect the others.

\highspace
\begin{examplebox}[: Logical Hadamard Gate]
    For many quantum error-correcting codes, a \textbf{logical Hadamard gate} can be implemented transversally by applying a Hadamard gate independently to every physical qubit:
    \begin{equation*}
        H_L = H^{\otimes n} = \underbrace{H \otimes H \otimes \ldots \otimes H}_{n \text{ times}}
    \end{equation*}
    Thus, the logical operation is realized through many independent physical operations performed in parallel.
\end{examplebox}

\highspace
\begin{examplebox}[: Logical CNOT Gate]
    A logical CNOT gate between two encoded logical qubits is often implemented by applying \textbf{physical CNOT gates pairwise between corresponding qubits of the two blocks}:
    \begin{center}
        \begin{quantikz}[row sep=0.25cm, column sep=0.7cm]
        \lstick{$A_1$} & \ctrl{1} & \\
        \lstick{$B_1$} & \targ{} & \\[0.2cm]

        \lstick{$A_2$} & \ctrl{1} & \\
        \lstick{$B_2$} & \targ{} & \\[0.2cm]

        \lstick{$A_3$} & \ctrl{1} & \\
        \lstick{$B_3$} & \targ{} & \\[0.2cm]

        \lstick{$A_4$} & \ctrl{1} & \\
        \lstick{$B_4$} & \targ{} &
    \end{quantikz}
    \end{center}
    Each pair of physical qubits interacts independently of the others.
\end{examplebox}